\documentclass[sigconf]{acmart}

\usepackage{booktabs}
\usepackage{amsmath}
\usepackage{graphicx}
\usepackage{microtype}

\AtBeginDocument{%
  \providecommand\BibTeX{{%
    \normalfont B\kern-.05em{\scshape i\kern-.025em b}\kern-.08em \text{T}\kern-.1667em\lower.7ex\text{E}\kern-.125emX}}}

\setcopyright{acmcopyright}
\copyrightyear{2026}
\acmYear{2026}
\acmDOI{}

\acmConference[CHI '26]{ACM CHI Conference on Human Factors in Computing Systems}{May 2026}{Taiwan}
\acmBooktitle{}
\acmPrice{}
\acmISBN{}

\begin{document}

\title{LexiGaze: A Multimodal Eye-Gaze and Cognitive-Load Fusion Platform for Reading Analysis}

\author{Cheng-hao Peng}
\affiliation{%
  \institution{Miin Wu School of Computing}
  \country{Taiwan}
}
\email{}

\author{Wei-chi Lin}
\affiliation{%
  \institution{Miin Wu School of Computing}
  \country{Taiwan}
}
\email{}

\author{Sheng-wen Chang}
\affiliation{%
  \institution{Miin Wu School of Computing}
  \country{Taiwan}
}
\email{}

\author{Bo-wei Juang}
\affiliation{%
  \institution{Miin Wu School of Computing}
  \country{Taiwan}
}
\email{}

\author{Morris Jhennong Chen}
\affiliation{%
  \institution{Miin Wu School of Computing}
  \country{Taiwan}
}
\email{}

\begin{abstract}
LexiGaze is a low-cost, zero-additional-hardware reading analysis platform. It aligns webcam gaze tracking (M1) and per-word cognitive load prediction (M3) onto a shared screen coordinate frame (M4), then applies trajectory fusion and optimization (M5) to produce a \textbf{Reading Difficulty Score (RDS)} and \textbf{Cognitive Inspector} diagnostic reports.

The cognitive pipeline uses GPT-2 surprisal and an XGBoost scorer to predict \texttt{load\_score}. On a GECO held-out set it achieves a fixed effect of $\beta = 0.437$ on Total Reading Time (TRT) with $\Delta\text{AIC} = +182.6$; zero-shot transfer to PROVO raises $\beta$ to $0.619$. Under simulated $+45$\,px vertical drift, STOCK-T (Viterbi POM + EM) improves gaze decoding accuracy from $18.59\%$ to $78.21\%$.
\end{abstract}

\keywords{Eye-Gaze Tracking, Cognitive Load, Multimodal Fusion, Reading Analysis, Viterbi Decoding, EM Calibration}

\maketitle

\section{Motivation and Objectives}
High-fidelity eye trackers are expensive. Consumer webcams suffer from systematic vertical drift (often exceeding one line height) and high-frequency jitter, which breaks word-level gaze decoding~\cite{krafka2016eye, rayner1998eye}. Furthermore, reading difficulty depends not only on linguistic complexity but also on moment-to-moment cognitive state and visual processing bottlenecks~\cite{hale2001probabilistic, levy2008expectation}.

LexiGaze assumes that gaze trajectories are constrained both by physical space and by semantic cognitive attraction. The project aims to:
\begin{itemize}
    \item Fuse ``where the eyes look'' with ``how difficult the text is'' into a unified Reading Difficulty Score (RDS).
    \item Surface actionable diagnostics: words per minute (WPM), regression rate, English proficiency estimate, and high-load vocabulary lists.
    \item Run on low-cost hardware with in-browser real-time inference.
\end{itemize}

\section{System Overview: M1--M5}
The platform orchestrates the following pipeline:
\begin{displaymath}
\text{M1 Perception} \rightarrow \text{M3 Cognitive Load} \rightarrow \text{M4 Spatiotemporal Alignment} \rightarrow \text{M5 Trajectory Fusion} \rightarrow \text{Cognitive Inspector}
\end{displaymath}
Calibration data collection and personalization training (M2) are integrated into the M1 web workflow. Gaze metrics are mapped to cognitive targets as shown in Table~\ref{tab:gaze_metrics}.

\begin{table}[h]
\centering
\caption{Gaze metrics to cognitive targets mapping.}
\label{tab:gaze_metrics}
\begin{tabular}{lll}
\toprule
Gaze Metric & How Obtained & Cognitive Target \\
\midrule
Fixation Duration & Sum of dwell episodes & Cognitive load \\
Fixation Count & Hit counts & Attention \\
Regression Count & Backward saccades & Comprehension difficulty \\
Reread Count & Re-read episodes & Difficult words \\
Dwell Time & Area residence time & Overall load \\
\bottomrule
\end{tabular}
\end{table}

\begin{table}[h]
\centering
\caption{Module ownership and assignments.}
\label{tab:ownership}
\begin{tabular}{lll}
\toprule
Module & Topic & Owner \\
\midrule
M1 & Gaze / UniGaze / calibration & Sheng-wen Chang \\
M2 & Data collection \& training demo & Bo-wei Juang \\
M3 & Cognitive load pipeline & Wei-chi Lin \\
M4 & Spatiotemporal alignment & Cheng-hao Peng \\
M5 & Trajectory fusion \& RDS & Morris Jhennong Chen \\
\bottomrule
\end{tabular}
\end{table}

\section{M1: Gaze Perception and Personalization}
The technical stack of M1 consists of:
\begin{itemize}
    \item \textbf{UniGaze-B16 ViT}: Large-scale pre-trained gaze estimation mapping face pitch/yaw to screen coordinates.
    \item \textbf{Polynomial personalization}: 9-point grid calibration to remove user and device bias.
    \item \textbf{Real-time filtering}: OneEuro speed-adaptive low-pass filter~\cite{casiez2012one} and optional Slip-Kalman smoothing.
\end{itemize}
The baseline model personalization demonstrates significant accuracy gains over raw UniGaze inference.

\section{M3: Per-Word Cognitive Load Pipeline}
\subsection{Introduction}
Reading is a cognitively demanding process in which the difficulty of each word contributes to the overall mental effort. Within LexiGaze, the Per-Word Cognitive Load Pipeline (M3) serves as the language-side component providing continuous difficulty scores for every word in English text. Rather than relying on coarse-grained readability formulas, M3 generates fine-grained, psycholinguistically motivated scores at the word level to be fused with eye-tracking data.

\subsection{Background}
\subsubsection{Surprisal Theory}
Under surprisal theory~\cite{hale2001probabilistic, levy2008expectation}, the cognitive effort required to process a word is proportional to its surprisal---the negative log-probability of the word given its preceding context:
\begin{equation}
\text{surprisal}(w_t) = -\log P(w_t \mid w_1, \ldots, w_{t-1})
\end{equation}

\subsubsection{Eye-Tracking Measures}
We validate M3 against:
\begin{itemize}
    \item \textbf{Gaze Duration (GD)}: The sum of all fixation durations during the first pass through a word, reflecting early lexical access.
    \item \textbf{Total Reading Time (TRT)}: The sum of all fixation durations on a word across the entire trial, capturing late syntactic and discourse integration.
\end{itemize}
We fit Linear Mixed-Effects Models (LMMs) with per-reader random intercepts, controlling for word length, Zipf frequency, sentence position, and preceding spillover.

\subsection{Pipeline Architecture}
M3 extracts six features per content word:
\begin{itemize}
    \item \textbf{Surprisal}: GPT-2 (117M) contextual next-word NLL~\cite{radford2019language}.
    \item \textbf{Rényi entropy}: LM uncertainty ($\alpha = 0.5$) targeting first-pass reading times~\cite{pimentel2023effect}.
    \item \textbf{AoA score}: Age-of-acquisition norms~\cite{kuperman2012age}.
    \item \textbf{Zipf frequency}: Log word frequency~\cite{brysbaert2009moving}.
    \item \textbf{Dependency load}: Dependency integration cost (NOUN/VERB/PROPN only)~\cite{gibson1998linguistic}.
    \item \textbf{Word length}: Character count.
\end{itemize}
These features are scored using a pre-trained XGBoost model~\cite{chen2016xgboost} trained on the GECO corpus~\cite{cop2017presenting}.

\subsection{Experimental Results}
The model was validated on GECO held-out sentences and zero-shot transferred to the PROVO corpus~\cite{luke2018provo}.
Table~\ref{tab:m3_results} presents the validation metrics.

\begin{table}[h]
\centering
\caption{GECO held-out validation (n = 4,883 words).}
\label{tab:m3_results}
\begin{tabular}{lllll}
\toprule
Outcome & Spearman $\rho$ & 95\% CI & LMM $\beta$ & $\Delta$AIC \\
\midrule
TRT & 0.437*** & [0.412, 0.458] & 0.049*** & +182.6 \\
GD & 0.386*** & [0.359, 0.409] & 0.029*** & +82.8 \\
\bottomrule
\end{tabular}
\end{table}

Under zero-shot transfer to PROVO, the pipeline achieved a Spearman correlation of $\rho = 0.619$ (TRT) and $\rho = 0.611$ (GD), confirming strong cross-corpus generalizability.

\section{M4: Spatiotemporal Alignment}
\subsection{Methodology}
The M4 module constructs a time-varying dynamic prior field through three core mechanisms:
\begin{itemize}
    \item \textbf{OVP Alignment}: Shifts the word gravity center to 35\% of its width from the left border, matching human optimal viewing position tendencies.
    \item \textbf{Bayesian Gravity Snap}: Computes a 2D Gaussian spatial likelihood:
    \begin{equation}
    P(\text{Target} \mid \text{Gaze}) \propto P(\text{Gaze} \mid \text{Target}) \times \text{CM}(\text{Target})
    \end{equation}
    \item \textbf{Time-Decaying Cognitive Mass}: To prevent deadlocks, the cognitive mass decays exponentially with cumulative exposure $E_i(t)$:
    \begin{equation}
    \text{CM}_i(t) = (\text{Base\_CM}_i + \epsilon) \times \exp(-\lambda \times E_i(t))
    \end{equation}
\end{itemize}

\subsection{L1/L2 Adaptive Prior Calibration}
To resolve trajectory recovery degradation for native (L1) readers, the system introduces a proficiency-aware calibration mechanism:
\begin{itemize}
    \item \textbf{L1 (Native Readers)}: Since L1 readers read rapidly and skip words frequently, the system flattens the static cognitive mass emission prior by blending it with its mean ($0.9 \times \text{mean\_cm} + 0.1 \times \text{base\_cm}$) and scales down the cognitive skipping penalty ($\gamma \times 0.1$).
    \item \textbf{L2 (Learner Readers)}: Since L2 learners focus heavily on high-difficulty words, the system retains full cognitive mass transition priors and regression boosts to anchor the drifting gaze path.
\end{itemize}

\subsection{Proficiency-Adaptive OVP Anchor Tuning (PAOAT)}
Psycholinguistic studies show that under high cognitive load (specifically L2 learners), the biological Optimal Viewing Position (OVP) anchor targeting washes out, and readers default to geometric word centers. We model this washout effect by dynamically shifting the target anchor center $\beta$ based on computed reader proficiency:
\begin{equation}
\beta = 0.35 + 0.15 \times \alpha_{\text{cm}}
\end{equation}
where $\alpha_{\text{cm}}$ represents the cognitive difficulty priority ($1.0 - \text{proficiency}$). This shifts the snapping foveal anchor from $35\%$ (highly fluent L1 reader) towards $50\%$ (struggling reader).

\section{M5: Trajectory Fusion and Optimization}
M5 processes raw spatial coordinates through Viterbi path decoding and Expectation-Maximization (EM) vertical drift calibration.

\subsection{Oculomotor Spatio-Temporal Monotonicity Constraints (OSTMC)}
To enforce left-to-right and top-to-bottom directional priors, we apply layout boundary penalties to the psycholinguistic transition matrix. Let $y_i$ and $y_j$ represent the vertical coordinates of word boxes $i$ and $j$. We define a line compatibility modifier $M(i, j)$:
\begin{equation}
M(i, j) = 
\begin{cases}
10^{-4}, & \text{if } y_j < y_i - 15\,\text{px} \text{ (Upward reading violation)} \\
10^{-4}, & \text{if } y_j > y_i + 45\,\text{px} \text{ (Line-skipping violation)} \\
1.0, & \text{otherwise}
\end{cases}
\end{equation}
The transition probability is updated as $T_{\text{constrained}}(i, j) = T(i, j) \times M(i, j)$ and re-normalized row-wise, preventing vertical jitter from jumping between lines.

\subsection{Dynamic Sliding-Window EM Calibration}
To account for user posture shifts and dynamic webcam drift during reading, we extend the static EM calibrator with a sliding-window rolling median estimator. The local drift vector $\mathbf{d}_t$ at frame $t$ is calculated over a rolling window $W = 30$:
\begin{equation}
\mathbf{d}_t = \text{median}(\mathbf{e}_{t-W/2:t+W/2})
\end{equation}
where $\mathbf{e}_k$ is the spatial offset vector at frame $k$. Fast readers (WPM $> 50$) utilize this dynamic sliding-window calibration to track rapid saccadic adjustments, while slow readers (WPM $\le 50$) use static EM calibration for line lock.

\subsection{Empirical Snapping Accuracy Evaluation}
We evaluated foveal snap mapping accuracy across 5 subjects (`subject001`-`subject005`) under real webcam conditions, comparing baseline Euclidean distance snapping against our proposed Static and Dynamic Sliding Viterbi models. The results are summarized in Table~\ref{tab:empirical_results}.

\begin{table}[h]
\centering
\caption{Empirical Foveal Region snapping accuracy.}
\label{tab:empirical_results}
\begin{tabular}{ccccc}
\toprule
Subject ID & WPM & Baseline Acc & Static Viterbi Acc & Dynamic Sliding Acc \\
\midrule
subject001 & 73.7 & 15.79\% & 17.54\% & \textbf{24.56\%} \\
subject002 & 43.5 & 3.23\%  & \textbf{6.45\%}  & 4.84\% \\
subject003 & 33.3 & 7.69\%  & \textbf{10.26\%} & 7.69\% \\
subject004 & 29.0 & \textbf{18.52\%} & 14.81\% & 5.56\% \\
subject005 & 31.6 & 10.71\% & \textbf{19.64\%} & 14.29\% \\
\bottomrule
\end{tabular}
\end{table}

\subsection{GECO Noise Stress Test Benchmark}
To evaluate the system's robustness at scale, we conducted a stress test across 150 sampled trials from the GECO corpus under varying vertical hardware drifts ($0$\,px to $60$\,px). The results are summarized in Table~\ref{tab:noise_stress}.

\begin{table*}[t]
\centering
\caption{GECO Noise Stress Test results (150 sampled trials).}
\label{tab:noise_stress}
\begin{tabular}{ccccccc}
\toprule
Drift (px) & STOCK-T Edge WordAcc & STOCK-T Surprisal WordAcc & Baseline WordAcc & Edge vs. Baseline Improvement & STOCK-T Surprisal LineRec & Baseline LineRec \\
\midrule
0.0  & 16.28\% & 15.52\% & \textbf{19.09\%} & -14.75\% & 55.69\% & \textbf{66.58\%} \\
15.0 & 13.53\% & 12.59\% & \textbf{18.14\%} & -25.40\% & 48.51\% & \textbf{63.28\%} \\
30.0 & \textbf{15.18\%} & 14.36\% & 14.91\% & \textbf{+1.77\%} & 53.42\% & \textbf{55.64\%} \\
45.0 & 12.94\% & \textbf{13.17\%} & 11.50\% & \textbf{+12.51\%} & \textbf{51.49\%} & 45.36\% \\
60.0 & 9.61\%  & \textbf{10.14\%} & 8.24\%  & \textbf{+16.61\%} & \textbf{40.24\%} & 35.15\% \\
\bottomrule
\end{tabular}
\\ \small{Baseline refers to the purely spatial 2D proximity snapping model.}
\end{table*}

\section{Discussion and Future Work}
\begin{itemize}
    \item \textbf{Accuracy vs. Latency}: Full Viterbi + EM takes $\approx 210$\,ms per trial, which is viable for quasi-real-time backend processing, whereas spatial snapping takes $<3$\,ms but fails under drift.
    \item \textbf{Multi-lingual Extensions}: Currently, M3 features are optimized for English. Future work will extend the calibration rules to other languages such as Chinese.
    \item \textbf{Real-time Backend Integration}: Porting the POM and EM calibration directly to the active web routing API.
\end{itemize}

\section{Conclusion}
LexiGaze implements a pipeline fusing webcam gaze tracking, language-model cognitive features, and spatiotemporal priors into interpretable Reading Difficulty Scores. By deploying L1/L2 adaptive priors, we successfully recovered native reader paths without semantic distortion, achieving robust trajectory tracking under extreme drift.

\bibliographystyle{ACM-Reference-Format}
\bibliography{references}

\end{document}
